\section{Actual monomial colex obstruction and intrinsic birth-block selections}
\label{sec:hironaka-monomial-birth-block-obstruction}

The following finite statements concern the actual ordered monomial
recipe. They do not assert an analytic example for the entire
history-sensitive canonical invariant. In particular, no identification
of that entire value with the concrete monomial data is assumed.

\begin{theorem}[An actual three-branch colex recipe changes under relabelling]
\label{hir:actual-monomial-three-branch-colex-permutation-obstruction}
\lean{Hironaka.HironakaMonomialBirthBlockObstruction.actual_colex_three_branch_obstruction}
For ambient bound three, three unit exponents and marking two, the
actual monomial invariant on the three indices is the key of the pair
$\{1,2\}$. Its value on $\{1,2\}$ is strictly greater than on
$\{0,1\}$. Swapping indices zero and two sends the selected pair to
$\{0,1\}$, a different pair. Thus this actual ordered recipe is not
invariant under arbitrary relabelling of equal initial branches.
\leanok
\uses{hir:monomial-invariant}
\end{theorem}
\begin{proof}
A subset of the three unit-weight indices is feasible exactly when it
has at least two members. The first key entry $3-|K|$ excludes the
three-member subset in favour of every feasible pair. Each feasible
pair has exponent sum two, so the last entry compares the pairs by
colexicographic order. Among the three pairs, $\{1,2\}$ is the unique
largest. On either two-member index set there is exactly one feasible
pair, and $\{0,1\}<_{\rm colex}\{1,2\}$. The indicated transposition
sends the larger pair to the smaller pair. These are evaluations of
the actual finite feasible-set and key definitions; no claim about
an unspecified full analytic invariant follows from them.
\end{proof}

\begin{theorem}[Invariance inside birth blocks is exactly whole-block selection]
\label{hir:actual-intrinsic-birth-block-selection-characterization}
\lean{Hironaka.HironakaMonomialBirthBlockObstruction.birthBlockClosed_iff_permutation_invariant}
Let a finite selected set of boundary labels have assigned birth or
colour labels. Its selection is invariant under every permutation
preserving these labels if and only if membership is constant inside
each birth or colour block. Equivalently, the selected set is a union
of entire such blocks.
\leanok
\uses{hir:actual-monomial-three-branch-colex-permutation-obstruction}
\end{theorem}
\begin{proof}
Suppose first that the selection is invariant. For two indices in the
same block, their transposition preserves every birth label: it swaps
two equal labels and fixes all others. Invariance of the finite set
under this transposition makes membership of the two indices
equivalent. Hence membership is constant on every block.

Conversely, suppose membership is constant on each block. A
birth-preserving permutation sends each selected index to another
index in its same block, which is therefore selected. For the reverse
inclusion, the inverse permutation also preserves birth labels; it
takes every selected target index to a selected preimage. Thus the
image of the actual finite selected set is exactly itself under each
allowed permutation. This uses the actual permutation action on the
indices, rather than a presumed invariance of the monomial recipe.
\end{proof}

\begin{corollary}[Three identical initial branches admit no symmetric minimal feasible pair]
\label{hir:actual-three-branch-no-symmetric-minimal-feasible-pair}
\lean{Hironaka.HironakaMonomialBirthBlockObstruction.no_permutation_invariant_two_of_three}
There is no two-element subset of three indistinguishable initial
branch labels which is invariant under all permutations. Consequently
a permutation-invariant replacement of the preceding actual monomial
recipe cannot retain its requirement of choosing one minimal feasible
pair on this input.
\leanok
\uses{hir:actual-intrinsic-birth-block-selection-characterization,hir:actual-monomial-three-branch-colex-permutation-obstruction}
\end{corollary}
\begin{proof}
Assign the same birth label to all three indices. The preceding
characterization says that an invariant selected set has constant
membership on all three. A two-element selected set is nonempty, so
choose one selected index. Every index must then be selected, forcing
the selected set to have cardinality three, a contradiction. For the
actual unit-exponent marking-two recipe, the minimal feasible subsets
are precisely the two-element subsets as already proved. The
obstruction requires a genuinely different intrinsic selection or
birth-block strategy; it supplies no strict-progress or termination
assertion for such a replacement.
\end{proof}
